\section{GPU Implementation}
\label{sec:gpu}
\sys{} separates RNS arithmetic, CKKS evaluation, neural-network operations, CUDA execution, and benchmark orchestration. Generic \texttt{HeOps} and \texttt{BootstrapOps} interfaces let the same trainer use CPU or device ciphertexts. The evaluator checks parameter identity, levels, scales, key availability, and device ownership before work. Counters commit after success; the CUDA bootstrap also waits for final device synchronization before merging its internal counters.

The CUDA backend uses device NTTs, basis conversion, pointwise arithmetic, hybrid key switching, and BSGS matvecs. Small public launch data use pinned transfer buffers. Reusing a hoisted decomposition avoids repeating the expensive preparation for related rotations. Device/CPU bit parity is meaningful only for cases actually executed; memory-skipped test functions can still report process success.

\paragraph{Streams and memory.}
Independent products may be issued to several streams, but a training step's dependent matvecs are executed sequentially. The measured pre-fix multi-stream benchmark was slower than sequential execution (Table~\ref{tab:gpu_vs_cpu_ops}). Calling-thread submission and pool handling were subsequently revised; their performance remeasurement is pending. An effective-bandwidth estimate for addition does not prove matvec occupancy or full-device saturation. That explanation remains a hypothesis until profiled.

The shared device has limited available memory. Training can upload keys capped to an explicitly selected level; this preserves the computation at that level, but does not measure the top-level workload. Allocation preflight uses estimated resident/key/workspace sizes. Pool trims occur after large setup stages, at epoch boundaries, and between ablation configurations. Estimated need and pool checkpoints are not a measured memory peak, and a successful preflight does not guarantee later allocations under contention.

\paragraph{Bootstrap boundary.}
The CUDA bootstrap implementation keeps ciphertext arithmetic on the device through SSE, modulus raising, transforms, and EvalMod. A narrow shared polynomial executor preserves the CPU operation order and rounding. Public integer constants are still materialized on the host and uploaded; future device timing must include that overhead. Local CPU replay and CUDA compilation support code readiness, not physical-GPU precision, latency, or memory feasibility. Those hardware measurements are pending.

The training benchmark still uses a CPU bootstrap adapter. Level operations run on the GPU; non-interactive refresh downloads to the CPU and uploads the result. A compiled device bootstrap API has not changed that benchmark. Its timing must therefore separate device work, transfers, and CPU refresh. Likewise interactive GPU training uses a client refresh. Neither path justifies an end-to-end all-GPU training claim.
